% Created 2026-03-08 Sun 15:33
% Intended LaTeX compiler: xelatex
\documentclass[12pt, letterpaper]{article}
\usepackage{graphicx}
\usepackage{longtable}
\usepackage{wrapfig}
\usepackage{rotating}
\usepackage[normalem]{ulem}
\usepackage{capt-of}
\usepackage{hyperref}
\usepackage[margin=1in]{geometry}
\usepackage{amsmath}
\usepackage{amssymb}
\usepackage{fontenc}
\usepackage{inputenc}
\usepackage{lmodern}
\usepackage{microtype}
\usepackage{booktabs}
\usepackage{array}
\usepackage{xcolor}
\usepackage{mdframed}
\usepackage{titlesec}
\usepackage{parskip}
\usepackage{hyperref}
\hypersetup{colorlinks=true, linkcolor=blue!50!black, urlcolor=blue!50!black}
\definecolor{medblue}{HTML}{2E4057}
\definecolor{accent}{HTML}{3A7CA5}
\definecolor{lightbg}{HTML}{E8F4F8}
\titleformat{\section}{\large\bfseries\color{medblue}}{\thesection}{1em}{}
\titleformat{\subsection}{\normalsize\bfseries\color{accent}}{\thesubsection}{1em}{}
\newmdenv[backgroundcolor=lightbg, linecolor=accent, linewidth=1.5pt, topline=true, bottomline=true, leftline=false, rightline=false, innerleftmargin=12pt, innerrightmargin=12pt, innertopmargin=10pt, innerbottommargin=10pt]{formulabox}
\newmdenv[backgroundcolor=white, linecolor=white, leftmargin=2cm, rightmargin=2cm, innerleftmargin=0pt, innerrightmargin=0pt, innertopmargin=4pt, innerbottommargin=4pt]{blockenv}
\author{Bradley K. DePuy, Independent Researcher}
\date{2026}
\title{The Medium — A Foundational Theory of Reality}
\hypersetup{
 pdfauthor={Bradley K. DePuy, Independent Researcher},
 pdftitle={The Medium — A Foundational Theory of Reality},
 pdfkeywords={},
 pdfsubject={},
 pdfcreator={Emacs 29.3 (Org mode 9.7.39)}, 
 pdflang={English}}
\begin{document}

\maketitle
\tableofcontents

\#\# -\textbf{- mode: org; coding: utf-8 -}-
\newpage
\begin{center}
\textbf{Version 0.3 (draft)  ·  DOI: \href{https://doi.org/10.5281/zenodo.18804587}{10.5281/zenodo.18804587}  ·  CC BY 4.0}
\end{center}

\begin{center}
\emph{Out of empty chaos all of creation becomes.}
\end{center}

That sentence is the theory. Everything that follows is the attempt to say it precisely.
\section{The Foundational Claim}
\label{sec:orge17c206}

The framework rests on a single foundational claim from which everything else follows:

\begin{quote}
The Medium always is and always was. It has no beginning and no end. It exists
outside of any temporal framework. Intrinsic distinction is an eternal property
of the Medium — not an event that occurred but a condition that always obtains.
The structure we observe is not the result of something that happened. It is the
expression of what the Medium always is.
\end{quote}

This claim does three things at once.

\textbf{It eliminates the need for a trigger.} The question of what initiated the
universe only arises if time is assumed to be primitive. If the Medium always is
and always was, there is no before and no initiation. The trigger problem
dissolves not by answering it but by identifying it as a consequence of a false
assumption.

\textbf{It eliminates time as a primitive.} The Medium exists outside of time. What
physics calls time is not a property of the substrate. It is an experience of
observers embedded within the differentiated structure of the Medium — a
consequence of transition, not a fundamental feature of reality.

\textbf{It reframes differentiation.} The emergence of structure from the Medium is not
a process that started. It is a condition that always obtains. The ground state
and the differentiated state coexist within the Medium eternally, because
intrinsic distinction — \(\iota\) — is eternal, and time does not apply to it.
\section{The Stochastic Character of the Medium}
\label{sec:org4907022}

Before the formal axioms are introduced, the nature of the Medium's ground state
requires precise characterization. That characterization is stochastic — and the
stochastic framing is not incidental. It is the most exact description available
of what the Medium is prior to the activation of intrinsic distinction.
\subsection{The Ground State as Fully Random Autonomous Stochastic System}
\label{sec:org97395c7}

The Medium at its ground state is a fully autonomous stochastic system. It
requires no inputs, no external cause, no temporal framework, and no governing
distribution. Its randomness is not randomness within a framework — it is
randomness without framework. No outcome is privileged over any other.
Non-existence holds no advantage over existence. Existence holds no advantage
over non-existence.

This is not chaos. Chaos retains structure. The ground state of the Medium
precedes structure entirely. It is the absolute suspension of all constraint,
including constraints upon randomness itself. It is meta-randomness — randomness
that is random about its own nature.
\subsection{What Fully Random Stochastic Nothingness Implies}
\label{sec:org73f509d}

In a fully random autonomous system of infinite extent with no preferred states,
every possible configuration has nonzero probability. This has a consequence
that is not intuitive but is mathematically unavoidable:

\begin{quote}
In an infinite substrate with fully random stochastic character and no
privileged states, every possible fluctuation — including the emergence of
ordered, structured, low-entropy configurations — is not merely possible but
necessary as a feature of the distribution itself.
\end{quote}

The universe is not something that happened against the odds. It is a necessary
expression of a substrate that has no mechanism for suppressing any possibility.
\subsection{The Relationship Between Stochastic Character and \(\iota\)}
\label{sec:org3b3fd02}

The stochastic character of the Medium describes the nature of the ground state.
Intrinsic distinction — \(\iota\) — describes the property that makes the ground
state generative rather than merely static. The relationship is precise:

\begin{itemize}
\item Stochastic character describes what the ground state \emph{is}
\item \(\iota\) describes what makes the ground state capable of producing structure
\item Manifestation \(V\) is the stochastic event — the actualization of a possibility
\item Structure formation at \(c\) is what a successful manifestation produces
\end{itemize}
\section{Axioms and Primitives}
\label{sec:org1c36b0e}

The formal definitions follow directly from the foundational claim and the
stochastic characterization of the ground state. Several are marked provisional.
Open questions are named precisely and carried explicitly as objects of
investigation.
\subsection{Axiom I — The Eternal Medium}
\label{sec:org7064d81}

The Medium always is and always was. It has no beginning, no end, and no
external cause.

\begin{formulabox}
\begin{equation}
\nexists\, t_0 : M \text{ begins at } t_0 \qquad \nexists\, t_1 : M \text{ ends at } t_1
\tag{1}
\end{equation}
\end{formulabox}

\emph{Status: Axiomatic.}
\subsection{Axiom II — Intrinsic Distinction and the Unity of V}
\label{sec:org3bfe750}

The Medium possesses intrinsic distinction, denoted \(\iota\). It is not an
operator, not a function, not a process. It is the one irreducible property of
the thing that simply is — the capacity of the undifferentiated to give rise to
the differentiated.

The activation of \(\iota\) at a location in a transition is denoted \(V\). \(V\)
is either present or absent. When present, it is identical at every location and
every transition — no magnitude, no internal structure, no degrees. A quantity
that is indivisible, identical at every occurrence, and admitting no degrees has
exactly one formal representation:

\begin{formulabox}
\begin{equation}
V \equiv 1
\tag{2}
\end{equation}
\centerline{\textit{V is dimensionless unity}}
\end{formulabox}

This is not \(V\) being assigned a value. It is \(V\) being recognized as what it
is. The Medium is pre-dimensional. No units exist at the substrate level.
\(V \equiv 1\) is a pure number — the simplest non-trivial element of the
substrate's arithmetic.

\emph{Status: Axiomatic. Nature of \(\iota\) under investigation.}
\subsection{Axiom III — Locality Emergence}
\label{sec:org7827361}

Location is not a primitive of the Medium. It does not exist prior to
differentiation. Location is produced when differentiation, propagating at the
limit \(c\), gives rise to self-sustaining structure. Before this condition is
met, the concepts of here and there, near and far, position and distance are
genuinely absent — not suppressed but nonexistent.

\emph{Status: Axiomatic.}
\subsection{Primitive I — The Medium (Post-Locality)}
\label{sec:orgad4a74d}

\begin{formulabox}
\begin{equation}
M = \{x \mid x \text{ is a location}\}
\tag{3}
\end{equation}
\end{formulabox}

\emph{Status: Provisional — post-locality approximation only.}
\subsection{Primitive II — The Population Function}
\label{sec:orga41ce0f}

With \(V \equiv 1\), the population function \(\phi\) maps each location to the
presence or absence of manifestation. The Medium is formally a binary field over
its locations.

\begin{formulabox}
\begin{align}
\phi &: M \to \{0, 1\} \tag{4} \\
\phi(x) &= 1 \text{ if manifestation has occurred at } x,\quad 0 \text{ otherwise} \tag{5}
\end{align}
\centerline{\textit{Both 0 and 1 are dimensionless. The substrate has no units.}}
\end{formulabox}

\emph{Status: Current.}
\subsection{Primitive III — The Manifestation Probability}
\label{sec:org2940d8d}

Manifestation is a binary stochastic event. At any transition \(t\), a location
\(x\) either manifests or it does not. The probability \(p(x)\) is a Bernoulli
parameter over the binary field.

\begin{formulabox}
\begin{align}
\text{manifestation}(x, t) &\in \{0, 1\} \tag{6} \\
p(x) &= \Pr(\phi(x) = 1), \qquad p : M \to [0,1] \tag{7}
\end{align}
\end{formulabox}

\emph{Status: Current.}
\subsection{Primitive IV — Transition}
\label{sec:orgfc2fe47}

Time is not a primitive of this framework. What physics describes as temporal
evolution is, at the level of the Medium, transition — a change from one binary
field state to another. The transition index \(t\) is a dimensionless label, not
a clock. No duration is implied.

\begin{formulabox}
\begin{equation}
M_t \to M_{t+1}
\tag{8}
\end{equation}
\end{formulabox}

\emph{Status: Current.}
\subsection{Primitive V — The Propagation Limit}
\label{sec:orga21b11d}

The Medium imposes a fundamental constraint on how rapidly the activation of
\(\iota\) can propagate across a transition. This limit is denoted \(c\). At the
substrate level \(c\) is dimensionless — a bound on the rate of change of the
binary field, not a velocity. What physics calls the speed of light is the
projection of this dimensionless bound into the observer's dimensional framework.

\begin{formulabox}
\begin{equation}
|\phi_{t+1}(x) - \phi_t(x)| \leq c \qquad \text{(dimensionless)}
\tag{9}
\end{equation}
\end{formulabox}

\emph{Status: Provisional.}
\subsection{Primitive VI — The Structure Formation Condition}
\label{sec:org40b8ff6}

A self-sustaining structure forms when propagation reaches the limit exactly.
This is the threshold at which locality is born.

\begin{formulabox}
\begin{equation}
|\phi_{t+1}(x) - \phi_t(x)| = c
\tag{10}
\end{equation}
\end{formulabox}

\emph{Status: Provisional.}
\section{Physics as Projection}
\label{sec:org011e58c}

The Medium is pre-dimensional. It has no meters, no seconds, no kilograms.
\(V \equiv 1\) is a pure number. The binary field \(\phi\) is dimensionless. The
propagation limit \(c\) is a dimensionless transition bound. The transition index
\(t\) is a dimensionless label. Nothing in the substrate description carries
units of any kind.

\textbf{Dimensions do not exist in the Medium. They are invented by observers.}

An observer is a self-sustaining structure in the binary field \(\phi\) — a
stable pattern that persists across transitions. Such a structure is embedded in
the differentiated region it constitutes. It cannot perceive the substrate
directly. It perceives only the relationships between structures like itself,
mediated by the propagation of manifestation events across the field. To describe
those relationships, it constructs a coordinate system — a framework of units and
dimensions that makes the relationships tractable.

That coordinate system is not a discovery about the substrate. It is a
construction of the observer. The dimensions of physics — length, time, mass,
charge, temperature — are categories the observer imposes on what it measures.
They have no counterpart in the substrate.

This has a precise consequence:

\begin{quote}
Every dimensional quantity in physics — every constant, every observable, every
equation — is a projection of a dimensionless substrate relationship onto the
observer's dimensional coordinate system. Physics, correctly understood from the
point of view of the Medium, is the complete and consistent description of those
projections. It is not wrong. It is a perfect map of the shadow. It is not a
description of the substrate that casts it.
\end{quote}
\subsection{The Projection Boundary}
\label{sec:org3ff47c1}

The boundary between the substrate and the observer's framework is precise. It
is the boundary between the dimensionless and the dimensional. Everything on the
substrate side is a pure number or a binary value. Everything on the observer
side carries units.

This boundary is a structural feature of the framework, not a gap in it. The
Medium framework identifies the boundary precisely and locates every quantity of
physics on the correct side of it.

\begin{table}[htbp]
\caption{\label{tab:projection}Substrate quantities and their observer projections. Each row is a derivation target.}
\centering
\begin{tabular}{lp{7cm}}
\toprule
\textbf{Substrate (dimensionless)} & \textbf{Observer projection (dimensional)}\\
\midrule
\(V \equiv 1\)  (pure number) & \(h\)  {[}action] — Planck's constant\\
\(c\)  (transition rate bound) & \(c\)  {[}length/time] — speed of light\\
\(\phi(x) \in \{0,1\}\)  (binary) & quantum amplitude, probability\\
\(t\)  (transition index) & time  [seconds]\\
structure formation condition & spacetime geometry\\
binary field entropy & thermodynamic entropy  [J/K]\\
propagation pattern density & energy, mass  [J, kg]\\
\(\iota\)  (intrinsic distinction) & {[}to be derived]\\
\bottomrule
\end{tabular}
\end{table}

The right column of Table \ref{tab:projection} is the language of physics. The left
column is the language of the Medium. The table is not a set of results. It is a
research program. Each row identifies a correspondence that must be formally
derived. Version 0.3 does not carry out these derivations. It names them
precisely and establishes the framework within which they can be pursued.
\subsection{What Projection Means for Planck's Constant}
\label{sec:org6441fb1}

The first row of Table \ref{tab:projection} is the most immediate. \(V \equiv 1\) at
the substrate level. Planck's constant \(h\) appears at the observer level as
the dimensional expression of that same event.

The relationship is not an equation between equals. \(V\) does not \emph{become} \(h\).
\(h\) is what an observer with a dimensional coordinate system measures when a
manifestation event \(V\) occurs.

\begin{formulabox}
\begin{align}
\text{Substrate:} \quad & V \equiv 1 \qquad \text{(dimensionless)} \tag{11} \\
\text{Observer:}  \quad & h = \mathrm{proj}(V) \qquad \text{(units of action, observer-dependent)} \tag{12}
\end{align}
\end{formulabox}

The numerical value of \(h\) in any given unit system encodes the relationship
between the observer's unit definitions and the substrate's natural scale. It is
not a fact about the substrate. It is a fact about the observer.

An open question is carried explicitly: why does the projection of \(V\) carry
units of action specifically? The answer requires a formal account of how the
observer's dimensional categories emerge from the structure of sustained binary
field configurations. That derivation is a primary target for future development.
\subsection{What Projection Means for the Speed of Light}
\label{sec:org54b6dbe}

The propagation limit \(c\) is dimensionless at the substrate level — a pure
bound on binary field transition rates. The observer, embedded in a structure
that extends across space and time, measures propagation in units of length per
unit time. The dimensionless bound is projected into this framework and appears
as a velocity.

Setting \(c = 1\) in physics — standard in relativistic calculations — is the
partial recovery of the substrate's natural description. It removes one layer of
the observer's unit construction. The Medium framework identifies why that
recovery works: the substrate never had the units to begin with.
\subsection{The General Principle}
\label{sec:org752be01}

The pattern is the same in every row of Table \ref{tab:projection}. A dimensionless
substrate relationship is measured by an observer whose coordinate system assigns
dimensions to everything it touches. The dimensional quantity that results is
real and useful — it correctly describes the structure of the observer's world.
But it carries information about the observer's measurement framework that is
absent from the substrate.

The presence of dimensional constants in the equations of physics — \(h\), \(c\),
\(G\), \(k\) — is the signature of this projection. Each constant marks a place
where the observer's dimensional framework has been imposed on a substrate
relationship that has none. In the substrate's own arithmetic, those constants do
not appear. The equations are clean.

The program implied by Table \ref{tab:projection} is the following: for each row,
derive the dimensional form of the observer quantity from the substrate quantity
and a formal account of the projection boundary. If that program succeeds, every
fundamental constant of physics will have been given a substrate derivation — and
the mystery of their values will have been resolved, not by finding deeper
constants, but by understanding that they were never fundamental to begin with.

\emph{Status: The projection boundary and its general principle are proposed as a
structural feature of the framework. The correspondence table is a research
program, not a set of results. Individual derivations are deferred to future
versions.}
\section{The Entropic Arc}
\label{sec:orgdf050ef}

The relationship between the stochastic ground state of the Medium and the
thermodynamic concept of entropy yields a consequence that bears directly on one
of the deepest unsolved problems in the current understanding of physical
reality.
\subsection{The Ground State as Maximum Entropy}
\label{sec:orgd962d53}

The ground state of the Medium — fully random, fully undifferentiated,
\(\phi = 0\) almost everywhere — is a maximum entropy state. Every configuration
of the binary field is equally probable.

\begin{formulabox}
\begin{equation}
S = k \log W \qquad \text{where } W \text{ is maximal at the ground state}
\tag{13}
\end{equation}
\centerline{\textit{Note: } $S$ \textit{ and } $k$ \textit{ are observer projections; } $W$ \textit{ is a dimensionless substrate count}}
\end{formulabox}
\subsection{Intrinsic Distinction as a Low Entropy Event}
\label{sec:org95be2ab}

The activation of \(\iota\) — a single location transitioning from \(\phi = 0\)
to \(\phi = 1\) — is necessarily a minimum local entropy event. One specific
binary configuration emerges from an infinite sea of equally probable
arrangements.
\subsection{Resolution of the Low Entropy Initial Condition Problem}
\label{sec:org66fc152}

The current understanding of physical reality faces a foundational problem: the
universe appears to have begun in a state of extraordinarily low entropy. Nothing
in the existing framework explains this.

The Medium framework resolves this without being asked. The first flip from 0 to
1 in the binary field is a low entropy event by definition — one specific outcome
from the maximum entropy ground state. The initial condition is not a brute fact.
It follows necessarily from Axiom II and the stochastic character of the ground
state.
\subsection{The Full Entropic Arc}
\label{sec:org526501e}

\begin{enumerate}
\item The Medium obtains eternally at maximum entropy — fully random, \(\phi = 0\) almost everywhere.
\item \(\iota\) manifests — \(V \equiv 1\) at a location — \(\phi\) flips from 0 to 1 — minimum local entropy — maximum information.
\item Structure forms at \(c\) — locality is born — binary patterns propagate and stabilize across transitions.
\item Complexity builds — sustained low-entropy configurations emerge as eddies of order in the binary field.
\item Consciousness appears — the Medium becoming locally aware of its own stochastic origin and entropic arc.
\item The trajectory continues toward maximum entropy — \(\phi\) returns to 0 almost everywhere.
\end{enumerate}

\begin{center}
\emph{Not a beginning and an end. A breathing.}
\end{center}

Because the Medium is eternal and \(\iota\) is its eternal property, this arc
does not occur once. It occurs wherever and whenever \(\iota\) manifests —
infinite stochastic fluctuations, infinite arcs, each spending its low-entropy
initial condition and returning to the ground. The binary field blinks into
structure and blinks out again, across an infinite substrate that does not
notice.
\section{The Problem With Where We Started}
\label{sec:orgd49e128}

Modern physics is extraordinary. Its achievements are beyond dispute. But the
foundations were built from the wrong starting point.

Human beings began making sense of the world from the human scale. That is
natural and inevitable. But it means the conceptual framework underlying physics
was constructed from the middle outward, not from the beginning. We built upward
toward the very large and downward toward the very small from a foundation that
was never examined as a foundation. It was simply where we happened to be
standing.

The consequences are visible at the edges of current theory. The two most
precisely verified theories in the history of science are fundamentally
incompatible with each other. The standard model requires unexplained phenomena
to account for observations that do not fit the prediction. Approximately 95\% of
the universe by current models consists of placeholder labels for things that do
not fit.

The presence of dimensional constants — \(h\), \(c\), \(G\) — as brute unexplained
inputs to the equations of physics is itself a symptom of this displaced
foundation. A theory built from the substrate would not require them. They appear
precisely because physics was built from the observer's position outward, and the
observer's dimensional coordinate system was baked into the language from the
start. The constants are the seams where the observer's framework was stitched
together without understanding what it was stitching.
\section{What the Framework Proposes}
\label{sec:org3591ebf}

\subsection{The Medium}
\label{sec:org1570299}

The Medium is the primitive substrate of reality — one undifferentiated thing
without parts, without position, without dimension. Its character is fully random
and autonomous. It is not a space containing things. It simply is, and \(\iota\)
is what it is.
\subsection{Nothing Moves}
\label{sec:org03d9251}

Nothing in the Medium moves. What we experience as motion is pattern — the
pattern of the binary field \(\phi\) shifting across the substrate across
transitions. What we observe as a particle is a sustained configuration in
\(\phi\) propagating at the dimensionless rate \(c\). Nothing traveled. The
pattern shifted.
\subsection{On Time}
\label{sec:org128ec23}

Time does not exist in the Medium framework as a fundamental property. What we
experience as the passage of time is the experience of transition — the change
from one binary field state to another, indexed by the dimensionless label \(t\).
The arrow of time is the arrow of increasing entropy in the binary field. Both
are projections of the observer's measurement framework onto substrate
transitions. They are real as experiences. They are not properties of the
substrate.
\subsection{On Existing Physics}
\label{sec:org9f2293c}

Once locality has been established through structure formation, the existing
mathematical frameworks of physics emerge as valid descriptions of differentiated
structure behavior within their respective domains. They are not wrong. They are
downstream. They are precise and consistent descriptions of the observer's
projection of the substrate. The Medium framework does not replace them. It
identifies the substrate from which they project — and defines the program by
which their dimensional constants can, in principle, be derived rather than
assumed.
\section{Unresolved Questions}
\label{sec:orga612033}

These are the frontier of the framework. Each is stated precisely and carried
explicitly.

\textbf{The Nature of Intrinsic Distinction.} What is \(\iota\)? The provisional
characterization — as the property permitting local, spontaneous entropy
reduction in a fully random maximum entropy binary substrate — is the central
open question of the theory.

\textbf{The Formal Account of the Projection Boundary.} How, precisely, does an
observer embedded in the binary field construct a dimensional coordinate system?
What determines which dimensional categories emerge? This is the foundational
question of Section 4 and the primary task for the next stage of formal
development.

\textbf{Why Does V Project to Action?} The observer's measurement of \(V \equiv 1\)
produces a quantity with units of action — energy multiplied by time. Why this
dimensional combination? The answer requires a derivation of the observer's
dimensional categories from the structure of sustained binary field
configurations. Until that derivation exists, the correspondence
\(V \equiv 1 \leftrightarrow h\) {[}action] is a named target, not a result.

\textbf{The Derivation of Each Row of Table 1.} Each correspondence in the projection
table is an open derivation: \(c\) {[}m/s] from the dimensionless transition bound;
\(G\) from structure formation geometry; thermodynamic entropy [J/K] from binary
field entropy. These are the specific formal tasks the framework now names and
defers.

\textbf{The Status of \(\alpha\).} The fine-structure constant \(\alpha \approx 1/137\)
is dimensionless. It cannot be a projection artifact in the same sense as \(h\)
and \(c\). Its status within the framework requires separate investigation.

\textbf{The Structure That Emerges.} What precisely is the self-sustaining binary
pattern that forms when the structure formation condition is met? The formal
definition of these patterns is the primary task of the next stage of
development.

\textbf{The Observer.} The observer is herself a sustained binary configuration within
\(\phi\), attempting to describe the substrate that produced her from within it.
Every formal system contains truths it cannot prove from within itself. The
Medium framework faces an analogous boundary, and acknowledges it.
\section{Summary of Axioms, Primitives, and Correspondences}
\label{sec:org0473528}

\begin{table}[htbp]
\caption{\label{tab:summary}Complete symbol table for version 0.3}
\centering
\begin{tabular}{llp{4cm}}
\toprule
\textbf{Symbol} & \textbf{Definition} & \textbf{Status}\\
\midrule
Stochastic ground & Fully random autonomous, no distribution & v0.2\\
\(\iota\) & Intrinsic Distinction — eternal, irreducible & Axiomatic\\
\(V \equiv 1\) & Manifestation event is dimensionless unity & New v0.3\\
\(\phi : M \to \{0,1\}\) & Population function — binary field, dimensionless & New v0.3\\
Projection boundary & Dimensional quantities are observer projections of substrate & New v0.3\\
\(h = \mathrm{proj}(V)\) & Planck's constant is the dimensional projection of \(V \equiv 1\) & New v0.3 · Prov.\\
\(c\) (substrate) & Propagation limit — dimensionless transition rate bound & New v0.3\\
\(c = \mathrm{proj}(c)\) & Speed of light is the dimensional projection of the transition bound & New v0.3 · Prov.\\
\(M\) & Eternal Medium — no beginning, no end & Axiomatic\\
\(p(x)\) & Manifestation probability — Bernoulli parameter & Current\\
\(t\) & Transition index — dimensionless label, not time & Current\\
\(S = k \log W\) & Boltzmann entropy — \(W\) is substrate count; \(S, k\) are projections & v0.2\\
\(\iota \Rightarrow S_{\min}\) & First distinction (\(\phi: 0 \to 1\)) is minimum local entropy & v0.2\\
\bottomrule
\end{tabular}
\end{table}
\section{A Personal Note}
\label{sec:orga28afff}

I am seventy-eight years old. I am not an academic. I have no credentials in
physics or mathematics. I began this work with a question that would not leave me
alone and a conviction that the question deserved a serious attempt at an answer.

I began this work knowing I may not live to see it reach its conclusions. That
does not trouble me. What matters is that the ideas are documented, attributed,
and made available to anyone who finds them worthy of development. If this
framework contains something true it belongs to whoever can take it further. My
name attached to it is enough — for the record, and for my family.

Gregor Mendel published his work on inheritance and it sat unread for decades
before being recognized as foundational. The ideas belong to whoever finds them.
The record belongs to their origin.
\section{License and Attribution}
\label{sec:org3a0a7c6}

This work is licensed under the Creative Commons Attribution 4.0 International
License (CC BY 4.0). You are free to share and adapt this material for any
purpose, including commercially, provided you give appropriate credit to the
original author, provide a link to the license, and indicate if changes were
made.

\url{https://creativecommons.org/licenses/by/4.0/}

All versions of this document are archived on Zenodo and remain linked to the
authorship record of Bradley K. DePuy, independent researcher.

\textbf{DOI:} \href{https://doi.org/10.5281/zenodo.18804587}{10.5281/zenodo.18804587}
\end{document}
